% ============================================================================
% Sección 2: Marco Teórico
% ============================================================================
\section{Marco Teórico}

\subsection{Fundamentos de búsqueda en grafos}

Un problema de búsqueda se define formalmente como una tupla $(S, s_0, G, A, c)$ donde $S$ es el conjunto de estados, $s_0 \in S$ es el estado inicial, $G \subseteq S$ es el conjunto de estados meta, $A$ es el conjunto de acciones disponibles, y $c: A \rightarrow \mathbb{R}^+$ es la función de costo \cite{russell2021artificial}. El objetivo es encontrar una secuencia de acciones que transforme $s_0$ en algún estado $g \in G$ minimizando el costo total.

\subsection{Algoritmos de búsqueda no informada}

\subsubsection{Búsqueda en Amplitud (BFS)}

BFS explora el espacio de estados nivel por nivel, expandiendo todos los nodos a profundidad $d$ antes de proceder a profundidad $d+1$. Emplea una cola FIFO como estructura de datos principal \cite{cormen2022introduction}.

\begin{itemize}
    \item \textbf{Completitud:} Sí, si el factor de ramificación $b$ es finito.
    \item \textbf{Optimalidad:} Sí, cuando todos los costos de aristas son uniformes.
    \item \textbf{Complejidad temporal:} $O(b^d)$
    \item \textbf{Complejidad espacial:} $O(b^d)$
\end{itemize}

\subsubsection{Búsqueda en Profundidad (DFS)}

DFS explora cada rama del árbol de búsqueda hasta su máxima profundidad antes de retroceder. Utiliza una pila (LIFO) o recursión como mecanismo de control \cite{cormen2022introduction}.

\begin{itemize}
    \item \textbf{Completitud:} No, en espacios infinitos sin detección de ciclos.
    \item \textbf{Optimalidad:} No garantizada.
    \item \textbf{Complejidad temporal:} $O(b^m)$ donde $m$ es la profundidad máxima.
    \item \textbf{Complejidad espacial:} $O(bm)$
\end{itemize}

\subsubsection{Algoritmo de Dijkstra}

El algoritmo de Dijkstra \cite{dijkstra1959note} generaliza BFS para grafos con costos no uniformes. Emplea una cola de prioridad ordenada por el costo acumulado $g(n)$ desde el origen.

\begin{itemize}
    \item \textbf{Completitud:} Sí, en grafos finitos con costos no negativos.
    \item \textbf{Optimalidad:} Sí, siempre.
    \item \textbf{Complejidad temporal:} $O((V + E) \log V)$ con cola de prioridad binaria.
    \item \textbf{Complejidad espacial:} $O(V)$
\end{itemize}

\subsection{Algoritmo A*}

\subsubsection{Definición formal}

El algoritmo A* \cite{hart1968formal} es un algoritmo de búsqueda informada que extiende el algoritmo de Dijkstra incorporando una función heurística. Para cada nodo $n$, A* evalúa:

\begin{equation}
    f(n) = g(n) + h(n)
    \label{eq:astar}
\end{equation}

donde:
\begin{itemize}
    \item $g(n)$: costo real del camino desde el nodo inicial hasta $n$.
    \item $h(n)$: estimación heurística del costo desde $n$ hasta la meta.
    \item $f(n)$: costo total estimado del camino óptimo a través de $n$.
\end{itemize}

\subsubsection{Pseudocódigo}

El Algoritmo~\ref{alg:astar} presenta el pseudocódigo formal de A*.

\begin{algorithm}[H]
\caption{Algoritmo A*}
\label{alg:astar}
\begin{algorithmic}[1]
\REQUIRE Problema de búsqueda $(s_0, G, \text{successors}, h)$
\ENSURE Camino óptimo de $s_0$ a $g \in G$ o fallo
\STATE $\text{open} \leftarrow$ cola de prioridad con $\{s_0\}$
\STATE $g[s_0] \leftarrow 0$
\STATE $\text{closed} \leftarrow \emptyset$
\WHILE{$\text{open} \neq \emptyset$}
    \STATE $n \leftarrow$ extraer nodo con menor $f(n)$ de open
    \IF{$n \in G$}
        \RETURN reconstruir\_camino($n$)
    \ENDIF
    \STATE $\text{closed} \leftarrow \text{closed} \cup \{n\}$
    \FORALL{$(m, c) \in \text{successors}(n)$}
        \IF{$m \in \text{closed}$}
            \STATE \textbf{continuar}
        \ENDIF
        \STATE $g_{\text{tentativo}} \leftarrow g[n] + c$
        \IF{$g_{\text{tentativo}} < g[m]$}
            \STATE $g[m] \leftarrow g_{\text{tentativo}}$
            \STATE $f[m] \leftarrow g[m] + h(m)$
            \STATE insertar o actualizar $m$ en open
        \ENDIF
    \ENDFOR
\ENDWHILE
\RETURN fallo
\end{algorithmic}
\end{algorithm}

\subsubsection{Propiedades formales}

\textbf{Definición 1 (Heurística admisible).} Una heurística $h$ es admisible si para todo nodo $n$:
\begin{equation}
    0 \leq h(n) \leq h^*(n)
    \label{eq:admissible}
\end{equation}
donde $h^*(n)$ es el costo real óptimo desde $n$ hasta la meta más cercana.

\textbf{Definición 2 (Heurística consistente).} Una heurística $h$ es consistente (o monótona) si para todo nodo $n$ y todo sucesor $m$ generado por la acción $a$:
\begin{equation}
    h(n) \leq c(n, a, m) + h(m)
    \label{eq:consistent}
\end{equation}

\textbf{Teorema 1 (Optimalidad de A*).} Si la heurística $h$ es admisible, A* con búsqueda en árbol encuentra la solución óptima. Si $h$ es consistente, A* con búsqueda en grafo encuentra la solución óptima \cite{hart1968formal}.

\textbf{Teorema 2 (Completitud).} A* es completo en grafos finitos o en grafos infinitos donde existe al menos una solución con costo finito y el costo mínimo de las aristas es positivo \cite{pearl1984heuristics}.

\textbf{Teorema 3 (Eficiencia óptima).} No existe otro algoritmo óptimo que, usando la misma heurística, garantice expandir menos nodos que A* (excluyendo empates en $f$) \cite{dechter1985generalized}.

\subsubsection{Complejidad}

\begin{itemize}
    \item \textbf{Temporal:} $O(b^d)$ en el peor caso, donde $b$ es el factor de ramificación y $d$ la profundidad de la solución. Con buena heurística, se reduce significativamente.
    \item \textbf{Espacial:} $O(b^d)$, ya que almacena todos los nodos generados.
\end{itemize}

El factor de ramificación efectivo $b^*$ permite evaluar la calidad de una heurística. Si A* genera $N$ nodos para encontrar una solución a profundidad $d$, entonces $b^*$ satisface:
\begin{equation}
    N + 1 = 1 + b^* + (b^*)^2 + \cdots + (b^*)^d
    \label{eq:branching}
\end{equation}

\subsection{Funciones heurísticas}

\subsubsection{Distancia Manhattan}

Para grids con movimiento en 4 direcciones (arriba, abajo, izquierda, derecha):
\begin{equation}
    h_{\text{Manhattan}}(n) = |x_n - x_g| + |y_n - y_g|
    \label{eq:manhattan}
\end{equation}

Es admisible para grids con costo uniforme y movimiento cardinal.

\subsubsection{Distancia Euclidiana}

Para grids con movimiento en 8 direcciones (incluye diagonales):
\begin{equation}
    h_{\text{Euclidiana}}(n) = \sqrt{(x_n - x_g)^2 + (y_n - y_g)^2}
    \label{eq:euclidean}
\end{equation}

Es admisible para cualquier esquema de movimiento, ya que la línea recta siempre subestima la distancia real del camino.

\subsection{Comparación teórica de algoritmos}

La Tabla~\ref{tab:comparison} sintetiza la comparación teórica entre los algoritmos analizados.

\begin{table}[H]
\centering
\caption{Comparación teórica de algoritmos de búsqueda}
\label{tab:comparison}
\renewcommand{\arraystretch}{1.3}
\begin{tabular}{@{}lcccc@{}}
\toprule
\textbf{Propiedad} & \textbf{BFS} & \textbf{DFS} & \textbf{Dijkstra} & \textbf{A*} \\
\midrule
Completo & Sí & No & Sí & Sí \\
Óptimo & Sí$^1$ & No & Sí & Sí$^2$ \\
Tiempo & $O(b^d)$ & $O(b^m)$ & $O((V\!+\!E)\log V)$ & $O(b^d)$ \\
Espacio & $O(b^d)$ & $O(bm)$ & $O(V)$ & $O(b^d)$ \\
Informado & No & No & No & Sí \\
Estructura & Cola & Pila & Heap & Heap \\
\bottomrule
\multicolumn{5}{l}{\footnotesize $^1$Con costos uniformes. $^2$Con heurística admisible.}
\end{tabular}
\end{table}
